\documentclass[aspectratio=169,11pt]{beamer}
\usetheme{Madrid}
\usecolortheme{default}
\setbeamertemplate{navigation symbols}{}
\setbeamertemplate{footline}[frame number]
\usepackage{enumitem}
\usepackage{amsmath,amssymb}
\usepackage{booktabs}
\usepackage{array}

\title{Databricks ML Associate}
\subtitle{Questions 701--800 (Eighth 100)}
\author{Practice Exam Deck}
\date{}

\begin{document}
	
	\begin{frame}
		\titlepage
	\end{frame}
	
	\begin{frame}{How to Use This Deck}
		\begin{itemize}
			\item 100 new questions (Q701--Q800), none repeated from Q1--Q700.
			\item Each question slide has four options (A--D).
			\item The next slide gives the correct answer and explains why the right option is right and why each wrong option is wrong.
			\item This set emphasizes Feature Store creation parameters, MLflow project/run tags, serving endpoint CRUD, registry operations, evaluator parameters, Hyperopt + MLflow, DLT modes, and CLI/API patterns.
		\end{itemize}
	\end{frame}
	
	% ---------------------------------------------------------------
	\begin{frame}{Question 701}
		Which parameter identifies the primary key when creating a Feature Store table?
		\begin{enumerate}[label=\Alph*.]
			\item \texttt{primary\_keys}
			\item \texttt{keys}
			\item \texttt{index\_cols}
			\item \texttt{row\_id}
		\end{enumerate}
	\end{frame}
	\begin{frame}{Answer 701}
		\textbf{Correct: A} --- \texttt{primary\_keys} defines the key columns.\\
		\vspace{0.5em}
		\textbf{Wrong:}
		\begin{itemize}
			\item B, C, D are not the Feature Store parameter.
		\end{itemize}
	\end{frame}
	
	% ---------------------------------------------------------------
	\begin{frame}{Question 702}
		Which parameter identifies the timestamp column when creating a time-series feature table?
		\begin{enumerate}[label=\Alph*.]
			\item \texttt{timestamp\_keys}
			\item \texttt{event\_time}
			\item \texttt{time\_col}
			\item \texttt{ts\_col}
		\end{enumerate}
	\end{frame}
	\begin{frame}{Answer 702}
		\textbf{Correct: A} --- \texttt{timestamp\_keys} defines time columns.\\
		\vspace{0.5em}
		\textbf{Wrong:}
		\begin{itemize}
			\item B, C, D are not the Feature Store parameter.
		\end{itemize}
	\end{frame}
	
	% ---------------------------------------------------------------
	\begin{frame}{Question 703}
		Which parameter gives a feature table a human-readable description?
		\begin{enumerate}[label=\Alph*.]
			\item \texttt{description}
			\item \texttt{comment}
			\item \texttt{metadata}
			\item \texttt{notes}
		\end{enumerate}
	\end{frame}
	\begin{frame}{Answer 703}
		\textbf{Correct: A} --- \texttt{description} sets the table description.\\
		\vspace{0.5em}
		\textbf{Wrong:}
		\begin{itemize}
			\item B, C, D are not the Feature Store parameter.
		\end{itemize}
	\end{frame}
	
	% ---------------------------------------------------------------
	\begin{frame}{Question 704}
		Which parameter specifies the database or schema for a feature table?
		\begin{enumerate}[label=\Alph*.]
			\item \texttt{database\_name} or schema name in the table name
			\item \texttt{catalog}
			\item \texttt{namespace}
			\item \texttt{prefix}
		\end{enumerate}
	\end{frame}
	\begin{frame}{Answer 704}
		\textbf{Correct: A} --- the table name includes the schema.\\
		\vspace{0.5em}
		\textbf{Wrong:}
		\begin{itemize}
			\item B, C, D are not the Feature Store parameter.
		\end{itemize}
	\end{frame}
	
	% ---------------------------------------------------------------
	\begin{frame}{Question 705}
		What does the \texttt{tags} parameter on a feature table do?
		\begin{enumerate}[label=\Alph*.]
			\item Attaches key-value metadata to the table
			\item Trains the model
			\item Serves the model
			\item Registers the model
		\end{enumerate}
	\end{frame}
	\begin{frame}{Answer 705}
		\textbf{Correct: A} --- tags attach metadata.\\
		\vspace{0.5em}
		\textbf{Wrong:}
		\begin{itemize}
			\item B, C, D are unrelated.
		\end{itemize}
	\end{frame}
	
	% ---------------------------------------------------------------
	\begin{frame}{Question 706}
		What is an MLflow project?
		\begin{enumerate}[label=\Alph*.]
			\item A packaged ML code project with an MLproject file
			\item A model
			\item A feature table
			\item A cluster
		\end{enumerate}
	\end{frame}
	\begin{frame}{Answer 706}
		\textbf{Correct: A} --- projects package code with an MLproject file.\\
		\vspace{0.5em}
		\textbf{Wrong:}
		\begin{itemize}
			\item B, C, D are unrelated.
		\end{itemize}
	\end{frame}
	
	% ---------------------------------------------------------------
	\begin{frame}{Question 707}
		What file defines an MLflow project?
		\begin{enumerate}[label=\Alph*.]
			\item \texttt{MLproject}
			\item \texttt{project.yaml}
			\item \texttt{mlflow.json}
			\item \texttt{requirements.txt}
		\end{enumerate}
	\end{frame}
	\begin{frame}{Answer 707}
		\textbf{Correct: A} --- \texttt{MLproject} defines a project.\\
		\vspace{0.5em}
		\textbf{Wrong:}
		\begin{itemize}
			\item B, C, D are not the project definition file.
		\end{itemize}
	\end{frame}
	
	% ---------------------------------------------------------------
	\begin{frame}{Question 708}
		What does an MLflow project's entry point define?
		\begin{enumerate}[label=\Alph*.]
			\item A runnable command and its parameters
			\item A model
			\item A feature table
			\item A cluster
		\end{enumerate}
	\end{frame}
	\begin{frame}{Answer 708}
		\textbf{Correct: A} --- entry points define runnable commands.\\
		\vspace{0.5em}
		\textbf{Wrong:}
		\begin{itemize}
			\item B, C, D are unrelated.
		\end{itemize}
	\end{frame}
	
	% ---------------------------------------------------------------
	\begin{frame}{Question 709}
		Which MLflow API runs an MLflow project?
		\begin{enumerate}[label=\Alph*.]
			\item \texttt{mlflow.run()}
			\item \texttt{mlflow.start()}
			\item \texttt{mlflow.execute()}
			\item \texttt{mlflow.launch()}
		\end{enumerate}
	\end{frame}
	\begin{frame}{Answer 709}
		\textbf{Correct: A} --- \texttt{mlflow.run()} runs a project.\\
		\vspace{0.5em}
		\textbf{Wrong:}
		\begin{itemize}
			\item B, C, D are not the API.
		\end{itemize}
	\end{frame}
	
	% ---------------------------------------------------------------
	\begin{frame}{Question 710}
		What is a run tag used for in MLflow?
		\begin{enumerate}[label=\Alph*.]
			\item Attaching metadata to a run
			\item Training a model
			\item Serving a model
			\item Registering a model
		\end{enumerate}
	\end{frame}
	\begin{frame}{Answer 710}
		\textbf{Correct: A} --- tags attach metadata.\\
		\vspace{0.5em}
		\textbf{Wrong:}
		\begin{itemize}
			\item B, C, D are unrelated.
		\end{itemize}
	\end{frame}
	
	% ---------------------------------------------------------------
	\begin{frame}{Question 711}
		Which MLflow API creates a serving endpoint?
		\begin{enumerate}[label=\Alph*.]
			\item \texttt{mlflow.deployments.create\_endpoint()}
			\item \texttt{mlflow.serve()}
			\item \texttt{mlflow.endpoint()}
			\item \texttt{mlflow.deploy()}
		\end{enumerate}
	\end{frame}
	\begin{frame}{Answer 711}
		\textbf{Correct: A} --- \texttt{create\_endpoint()} creates endpoints.\\
		\vspace{0.5em}
		\textbf{Wrong:}
		\begin{itemize}
			\item B, C, D are not the API.
		\end{itemize}
	\end{frame}
	
	% ---------------------------------------------------------------
	\begin{frame}{Question 712}
		Which MLflow API queries a serving endpoint?
		\begin{enumerate}[label=\Alph*.]
			\item \texttt{mlflow.deployments.predict()}
			\item \texttt{mlflow.predict()}
			\item \texttt{mlflow.query()}
			\item \texttt{mlflow.infer()}
		\end{enumerate}
	\end{frame}
	\begin{frame}{Answer 712}
		\textbf{Correct: A} --- \texttt{predict()} queries endpoints.\\
		\vspace{0.5em}
		\textbf{Wrong:}
		\begin{itemize}
			\item B, C, D are not the API.
		\end{itemize}
	\end{frame}
	
	% ---------------------------------------------------------------
	\begin{frame}{Question 713}
		Which MLflow API lists serving endpoints?
		\begin{enumerate}[label=\Alph*.]
			\item \texttt{mlflow.deployments.list\_endpoints()}
			\item \texttt{mlflow.endpoints()}
			\item \texttt{mlflow.list\_serve()}
			\item \texttt{mlflow.deploy\_list()}
		\end{enumerate}
	\end{frame}
	\begin{frame}{Answer 713}
		\textbf{Correct: A} --- \texttt{list\_endpoints()} lists endpoints.\\
		\vspace{0.5em}
		\textbf{Wrong:}
		\begin{itemize}
			\item B, C, D are not the API.
		\end{itemize}
	\end{frame}
	
	% ---------------------------------------------------------------
	\begin{frame}{Question 714}
		Which MLflow API deletes a serving endpoint?
		\begin{enumerate}[label=\Alph*.]
			\item \texttt{mlflow.deployments.delete\_endpoint()}
			\item \texttt{mlflow.delete\_serve()}
			\item \texttt{mlflow.remove\_endpoint()}
			\item \texttt{mlflow.endpoint\_delete()}
		\end{enumerate}
	\end{frame}
	\begin{frame}{Answer 714}
		\textbf{Correct: A} --- \texttt{delete\_endpoint()} deletes endpoints.\\
		\vspace{0.5em}
		\textbf{Wrong:}
		\begin{itemize}
			\item B, C, D are not the API.
		\end{itemize}
	\end{frame}
	
	% ---------------------------------------------------------------
	\begin{frame}{Question 715}
		What does a serving endpoint's traffic configuration control?
		\begin{enumerate}[label=\Alph*.]
			\item How requests are routed to served entities
			\item The training algorithm
			\item The feature table
			\item The cluster size
		\end{enumerate}
	\end{frame}
	\begin{frame}{Answer 715}
		\textbf{Correct: A} --- traffic configs route requests.\\
		\vspace{0.5em}
		\textbf{Wrong:}
		\begin{itemize}
			\item B, C, D are unrelated.
		\end{itemize}
	\end{frame}
	
	% ---------------------------------------------------------------
	\begin{frame}{Question 716}
		What is a served entity in a serving endpoint?
		\begin{enumerate}[label=\Alph*.]
			\item A model version or feature spec the endpoint serves
			\item A training job
			\item A feature table
			\item A cluster
		\end{enumerate}
	\end{frame}
	\begin{frame}{Answer 716}
		\textbf{Correct: A} --- served entities define what is served.\\
		\vspace{0.5em}
		\textbf{Wrong:}
		\begin{itemize}
			\item B, C, D are unrelated.
		\end{itemize}
	\end{frame}
	
	% ---------------------------------------------------------------
	\begin{frame}{Question 717}
		Which MLflow API creates a registered model in Unity Catalog?
		\begin{enumerate}[label=\Alph*.]
			\item \texttt{mlflow.create\_model()} or \texttt{register\_model()}
			\item \texttt{mlflow.new\_model()}
			\item \texttt{mlflow.add\_model()}
			\item \texttt{mlflow.make\_model()}
		\end{enumerate}
	\end{frame}
	\begin{frame}{Answer 717}
		\textbf{Correct: A} --- models are created via those APIs.\\
		\vspace{0.5em}
		\textbf{Wrong:}
		\begin{itemize}
			\item B, C, D are not the API.
		\end{itemize}
	\end{frame}
	
	% ---------------------------------------------------------------
	\begin{frame}{Question 718}
		Which MLflow API sets a registered model's description?
		\begin{enumerate}[label=\Alph*.]
			\item \texttt{mlflow.update\_model()}
			\item \texttt{mlflow.set\_model\_description()}
			\item \texttt{mlflow.describe\_model()}
			\item \texttt{mlflow.model\_info()}
		\end{enumerate}
	\end{frame}
	\begin{frame}{Answer 718}
		\textbf{Correct: A} --- \texttt{update\_model()} updates model metadata.\\
		\vspace{0.5em}
		\textbf{Wrong:}
		\begin{itemize}
			\item B, C, D are not the API.
		\end{itemize}
	\end{frame}
	
	% ---------------------------------------------------------------
	\begin{frame}{Question 719}
		Which MLflow API transitions a model version's stage?
		\begin{enumerate}[label=\Alph*.]
			\item \texttt{mlflow.transition\_model\_version\_stage()}
			\item \texttt{mlflow.set\_stage()}
			\item \texttt{mlflow.move\_version()}
			\item \texttt{mlflow.change\_stage()}
		\end{enumerate}
	\end{frame}
	\begin{frame}{Answer 719}
		\textbf{Correct: A} --- that API transitions stages.\\
		\vspace{0.5em}
		\textbf{Wrong:}
		\begin{itemize}
			\item B, C, D are not the API.
		\end{itemize}
	\end{frame}
	
	% ---------------------------------------------------------------
	\begin{frame}{Question 720}
		Which MLflow API deletes a model version?
		\begin{enumerate}[label=\Alph*.]
			\item \texttt{mlflow.delete\_model\_version()}
			\item \texttt{mlflow.remove\_version()}
			\item \texttt{mlflow.drop\_version()}
			\item \texttt{mlflow.model\_delete()}
		\end{enumerate}
	\end{frame}
	\begin{frame}{Answer 720}
		\textbf{Correct: A} --- \texttt{delete\_model\_version()} deletes versions.\\
		\vspace{0.5em}
		\textbf{Wrong:}
		\begin{itemize}
			\item B, C, D are not the API.
		\end{itemize}
	\end{frame}
	
	% ---------------------------------------------------------------
	\begin{frame}{Question 721}
		Which MLflow API updates a model version's description?
		\begin{enumerate}[label=\Alph*.]
			\item \texttt{mlflow.update\_model\_version()}
			\item \texttt{mlflow.set\_version\_desc()}
			\item \texttt{mlflow.describe\_version()}
			\item \texttt{mlflow.version\_info()}
		\end{enumerate}
	\end{frame}
	\begin{frame}{Answer 721}
		\textbf{Correct: A} --- \texttt{update\_model\_version()} updates metadata.\\
		\vspace{0.5em}
		\textbf{Wrong:}
		\begin{itemize}
			\item B, C, D are not the API.
		\end{itemize}
	\end{frame}
	
	% ---------------------------------------------------------------
	\begin{frame}{Question 722}
		What does \texttt{setLabelCol()} do in a Spark ML evaluator?
		\begin{enumerate}[label=\Alph*.]
			\item Sets the target column for evaluation
			\item Sets the metric
			\item Sets the prediction column
			\item Sets the model
		\end{enumerate}
	\end{frame}
	\begin{frame}{Answer 722}
		\textbf{Correct: A} --- it sets the label column.\\
		\vspace{0.5em}
		\textbf{Wrong:}
		\begin{itemize}
			\item B is \texttt{setMetricName}.
			\item C is \texttt{setPredictionCol}.
			\item D is unrelated.
		\end{itemize}
	\end{frame}
	
	% ---------------------------------------------------------------
	\begin{frame}{Question 723}
		What does \texttt{setRawPredictionCol()} do in a Spark ML evaluator?
		\begin{enumerate}[label=\Alph*.]
			\item Sets the column containing raw prediction scores
			\item Sets the target column
			\item Sets the metric
			\item Sets the model
		\end{enumerate}
	\end{frame}
	\begin{frame}{Answer 723}
		\textbf{Correct: A} --- it sets the raw prediction column.\\
		\vspace{0.5em}
		\textbf{Wrong:}
		\begin{itemize}
			\item B is \texttt{setLabelCol}.
			\item C is \texttt{setMetricName}.
			\item D is unrelated.
		\end{itemize}
	\end{frame}
	
	% ---------------------------------------------------------------
	\begin{frame}{Question 724}
		What does \texttt{setProbabilityCol()} do in a Spark ML evaluator?
		\begin{enumerate}[label=\Alph*.]
			\item Sets the column containing predicted probabilities
			\item Sets the target column
			\item Sets the metric
			\item Sets the model
		\end{enumerate}
	\end{frame}
	\begin{frame}{Answer 724}
		\textbf{Correct: A} --- it sets the probability column.\\
		\vspace{0.5em}
		\textbf{Wrong:}
		\begin{itemize}
			\item B is \texttt{setLabelCol}.
			\item C is \texttt{setMetricName}.
			\item D is unrelated.
		\end{itemize}
	\end{frame}
	
	% ---------------------------------------------------------------
	\begin{frame}{Question 725}
		Which Spark ML evaluator computes \texttt{areaUnderPR}?
		\begin{enumerate}[label=\Alph*.]
			\item \texttt{BinaryClassificationEvaluator}
			\item \texttt{RegressionEvaluator}
			\item \texttt{ClusteringEvaluator}
			\item \texttt{MulticlassClassificationEvaluator}
		\end{enumerate}
	\end{frame}
	\begin{frame}{Answer 725}
		\textbf{Correct: A} --- binary evaluator computes \texttt{areaUnderPR}.\\
		\vspace{0.5em}
		\textbf{Wrong:}
		\begin{itemize}
			\item B, C, D are other evaluators.
		\end{itemize}
	\end{frame}
	
	% ---------------------------------------------------------------
	\begin{frame}{Question 726}
		Which Spark ML evaluator supports \texttt{accuracy}?
		\begin{enumerate}[label=\Alph*.]
			\item \texttt{MulticlassClassificationEvaluator}
			\item \texttt{RegressionEvaluator}
			\item \texttt{ClusteringEvaluator}
			\item \texttt{BinaryClassificationEvaluator}
		\end{enumerate}
	\end{frame}
	\begin{frame}{Answer 726}
		\textbf{Correct: A} --- multiclass evaluator supports accuracy.\\
		\vspace{0.5em}
		\textbf{Wrong:}
		\begin{itemize}
			\item B, C, D are other evaluators.
		\end{itemize}
	\end{frame}
	
	% ---------------------------------------------------------------
	\begin{frame}{Question 727}
		Which Spark ML evaluator supports \texttt{mae}?
		\begin{enumerate}[label=\Alph*.]
			\item \texttt{RegressionEvaluator}
			\item \texttt{BinaryClassificationEvaluator}
			\item \texttt{MulticlassClassificationEvaluator}
			\item \texttt{ClusteringEvaluator}
		\end{enumerate}
	\end{frame}
	\begin{frame}{Answer 727}
		\textbf{Correct: A} --- regression evaluator supports MAE.\\
		\vspace{0.5em}
		\textbf{Wrong:}
		\begin{itemize}
			\item B, C, D are other evaluators.
		\end{itemize}
	\end{frame}
	
	% ---------------------------------------------------------------
	\begin{frame}{Question 728}
		What is an MLflow autologging benefit?
		\begin{enumerate}[label=\Alph*.]
			\item It removes the need to log params/metrics manually
			\item It trains faster
			\item It serves faster
			\item It registers models
		\end{enumerate}
	\end{frame}
	\begin{frame}{Answer 728}
		\textbf{Correct: A} --- autologging automates logging.\\
		\vspace{0.5em}
		\textbf{Wrong:}
		\begin{itemize}
			\item B, C, D are false.
		\end{itemize}
	\end{frame}
	
	% ---------------------------------------------------------------
	\begin{frame}{Question 729}
		How does Hyperopt integrate with MLflow?
		\begin{enumerate}[label=\Alph*.]
			\item By logging each trial's results to MLflow when enabled
			\item By replacing MLflow
			\item By serving models
			\item By registering models
		\end{enumerate}
	\end{frame}
	\begin{frame}{Answer 729}
		\textbf{Correct: A} --- Hyperopt can log trials to MLflow.\\
		\vspace{0.5em}
		\textbf{Wrong:}
		\begin{itemize}
			\item B, C, D are false.
		\end{itemize}
	\end{frame}
	
	% ---------------------------------------------------------------
	\begin{frame}{Question 730}
		Which Hyperopt feature logs trial results to MLflow?
		\begin{enumerate}[label=\Alph*.]
			\item \texttt{MLflowLogger} or auto-logging integration
			\item \texttt{SparkTrials} only
			\item \texttt{STATUS\_OK}
			\item \texttt{fmin}
		\end{enumerate}
	\end{frame}
	\begin{frame}{Answer 730}
		\textbf{Correct: A} --- MLflow logging is enabled via the integration.\\
		\vspace{0.5em}
		\textbf{Wrong:}
		\begin{itemize}
			\item B is distribution.
			\item C is a status.
			\item D is the optimizer.
		\end{itemize}
	\end{frame}
	
	% ---------------------------------------------------------------
	\begin{frame}{Question 731}
		What is a Delta Live Tables pipeline's continuous mode?
		\begin{enumerate}[label=\Alph*.]
			\item It keeps the pipeline running to process new data
			\item It runs once and stops
			\item It trains models
			\item It serves models
		\end{enumerate}
	\end{frame}
	\begin{frame}{Answer 731}
		\textbf{Correct: A} --- continuous pipelines keep processing new data.\\
		\vspace{0.5em}
		\textbf{Wrong:}
		\begin{itemize}
			\item B is triggered mode.
			\item C, D are unrelated.
		\end{itemize}
	\end{frame}
	
	% ---------------------------------------------------------------
	\begin{frame}{Question 732}
		What is a Delta Live Tables pipeline's triggered mode?
		\begin{enumerate}[label=\Alph*.]
			\item It processes data and then stops
			\item It runs continuously
			\item It trains models
			\item It serves models
		\end{enumerate}
	\end{frame}
	\begin{frame}{Answer 732}
		\textbf{Correct: A} --- triggered pipelines run and stop.\\
		\vspace{0.5em}
		\textbf{Wrong:}
		\begin{itemize}
			\item B is continuous.
			\item C, D are unrelated.
		\end{itemize}
	\end{frame}
	
	% ---------------------------------------------------------------
	\begin{frame}{Question 733}
		What is a Delta Live Tables pipeline's event log used for?
		\begin{enumerate}[label=\Alph*.]
			\item Tracking pipeline execution and data quality
			\item Training models
			\item Serving models
			\item Registering models
		\end{enumerate}
	\end{frame}
	\begin{frame}{Answer 733}
		\textbf{Correct: A} --- event logs track pipeline execution.\\
		\vspace{0.5em}
		\textbf{Wrong:}
		\begin{itemize}
			\item B, C, D are unrelated.
		\end{itemize}
	\end{frame}
	
	% ---------------------------------------------------------------
	\begin{frame}{Question 734}
		What is a Delta Live Tables pipeline's materialized view?
		\begin{enumerate}[label=\Alph*.]
			\item A precomputed query result stored as a table
			\item A model
			\item A feature table
			\item A cluster
		\end{enumerate}
	\end{frame}
	\begin{frame}{Answer 734}
		\textbf{Correct: A} --- materialized views store query results.\\
		\vspace{0.5em}
		\textbf{Wrong:}
		\begin{itemize}
			\item B, C, D are unrelated.
		\end{itemize}
	\end{frame}
	
	% ---------------------------------------------------------------
	\begin{frame}{Question 735}
		Which is a valid Databricks CLI command to list jobs?
		\begin{enumerate}[label=\Alph*.]
			\item \texttt{databricks jobs list}
			\item \texttt{databricks list jobs}
			\item \texttt{databricks jobs all}
			\item \texttt{databricks job ls}
		\end{enumerate}
	\end{frame}
	\begin{frame}{Answer 735}
		\textbf{Correct: A} --- \texttt{jobs list} lists jobs.\\
		\vspace{0.5em}
		\textbf{Wrong:}
		\begin{itemize}
			\item B, C, D are not the command syntax.
		\end{itemize}
	\end{frame}
	
	% ---------------------------------------------------------------
	\begin{frame}{Question 736}
		Which is a valid Databricks CLI command to run a job?
		\begin{enumerate}[label=\Alph*.]
			\item \texttt{databricks jobs run-now}
			\item \texttt{databricks run job}
			\item \texttt{databricks jobs start}
			\item \texttt{databricks job execute}
		\end{enumerate}
	\end{frame}
	\begin{frame}{Answer 736}
		\textbf{Correct: A} --- \texttt{jobs run-now} starts a job.\\
		\vspace{0.5em}
		\textbf{Wrong:}
		\begin{itemize}
			\item B, C, D are not the CLI syntax.
		\end{itemize}
	\end{frame}
	
	% ---------------------------------------------------------------
	\begin{frame}{Question 737}
		Which is a valid Databricks CLI command to list clusters?
		\begin{enumerate}[label=\Alph*.]
			\item \texttt{databricks clusters list}
			\item \texttt{databricks list clusters}
			\item \texttt{databricks clusters all}
			\item \texttt{databricks cluster ls}
		\end{enumerate}
	\end{frame}
	\begin{frame}{Answer 737}
		\textbf{Correct: A} --- \texttt{clusters list} lists clusters.\\
		\vspace{0.5em}
		\textbf{Wrong:}
		\begin{itemize}
			\item B, C, D are not the CLI syntax.
		\end{itemize}
	\end{frame}
	
	% ---------------------------------------------------------------
	\begin{frame}{Question 738}
		Which REST API endpoint lists jobs?
		\begin{enumerate}[label=\Alph*.]
			\item \texttt{GET /api/2.1/jobs/list}
			\item \texttt{GET /api/2.1/jobs/all}
			\item \texttt{GET /api/2.1/list/jobs}
			\item \texttt{GET /api/2.1/jobs}
		\end{enumerate}
	\end{frame}
	\begin{frame}{Answer 738}
		\textbf{Correct: A} --- \texttt{/jobs/list} lists jobs.\\
		\vspace{0.5em}
		\textbf{Wrong:}
		\begin{itemize}
			\item B, C, D are not the correct endpoint.
		\end{itemize}
	\end{frame}
	
	% ---------------------------------------------------------------
	\begin{frame}{Question 739}
		Which REST API endpoint triggers a job run?
		\begin{enumerate}[label=\Alph*.]
			\item \texttt{POST /api/2.1/jobs/run-now}
			\item \texttt{POST /api/2.1/jobs/start}
			\item \texttt{POST /api/2.1/jobs/execute}
			\item \texttt{POST /api/2.1/jobs/launch}
		\end{enumerate}
	\end{frame}
	\begin{frame}{Answer 739}
		\textbf{Correct: A} --- \texttt{run-now} triggers a run.\\
		\vspace{0.5em}
		\textbf{Wrong:}
		\begin{itemize}
			\item B, C, D are not the endpoint.
		\end{itemize}
	\end{frame}
	
	% ---------------------------------------------------------------
	\begin{frame}{Question 740}
		Which REST API endpoint lists clusters?
		\begin{enumerate}[label=\Alph*.]
			\item \texttt{GET /api/2.0/clusters/list}
			\item \texttt{GET /api/2.0/clusters/all}
			\item \texttt{GET /api/2.0/clusters}
			\item \texttt{GET /api/2.0/list/clusters}
		\end{enumerate}
	\end{frame}
	\begin{frame}{Answer 740}
		\textbf{Correct: A} --- \texttt{/clusters/list} lists clusters.\\
		\vspace{0.5em}
		\textbf{Wrong:}
		\begin{itemize}
			\item B, C, D are not the endpoint.
		\end{itemize}
	\end{frame}
	
	% ---------------------------------------------------------------
	\begin{frame}{Question 741}
		What is a Databricks personal access token used for?
		\begin{enumerate}[label=\Alph*.]
			\item Authenticating API and CLI calls
			\item Training models
			\item Serving models
			\item Registering models
		\end{enumerate}
	\end{frame}
	\begin{frame}{Answer 741}
		\textbf{Correct: A} --- PATs authenticate API/CLI calls.\\
		\vspace{0.5em}
		\textbf{Wrong:}
		\begin{itemize}
			\item B, C, D are unrelated.
		\end{itemize}
	\end{frame}
	
	% ---------------------------------------------------------------
	\begin{frame}{Question 742}
		Which is a best practice for personal access tokens?
		\begin{enumerate}[label=\Alph*.]
			\item Rotate tokens regularly and store them securely
			\item Share tokens publicly
			\item Commit tokens to Git
			\item Use the same token for all users
		\end{enumerate}
	\end{frame}
	\begin{frame}{Answer 742}
		\textbf{Correct: A} --- rotate and store tokens securely.\\
		\vspace{0.5em}
		\textbf{Wrong:}
		\begin{itemize}
			\item B, C, D are unsafe.
		\end{itemize}
	\end{frame}
	
	% ---------------------------------------------------------------
	\begin{frame}{Question 743}
		What is OAuth used for in Databricks?
		\begin{enumerate}[label=\Alph*.]
			\item Delegated authentication for users and service principals
			\item Training models
			\item Serving models
			\item Registering models
		\end{enumerate}
	\end{frame}
	\begin{frame}{Answer 743}
		\textbf{Correct: A} --- OAuth supports delegated authentication.\\
		\vspace{0.5em}
		\textbf{Wrong:}
		\begin{itemize}
			\item B, C, D are unrelated.
		\end{itemize}
	\end{frame}
	
	% ---------------------------------------------------------------
	\begin{frame}{Question 744}
		What is a Unity Catalog model's fully qualified name format?
		\begin{enumerate}[label=\Alph*.]
			\item \texttt{catalog.schema.model}
			\item \texttt{database.model}
			\item \texttt{workspace.model}
			\item \texttt{hive.model}
		\end{enumerate}
	\end{frame}
	\begin{frame}{Answer 744}
		\textbf{Correct: A} --- UC models use \texttt{catalog.schema.model}.\\
		\vspace{0.5em}
		\textbf{Wrong:}
		\begin{itemize}
			\item B, C, D are not UC formats.
		\end{itemize}
	\end{frame}
	
	% ---------------------------------------------------------------
	\begin{frame}{Question 745}
		Which privilege is needed to register a model in Unity Catalog?
		\begin{enumerate}[label=\Alph*.]
			\item \texttt{USE CATALOG} and \texttt{USE SCHEMA} plus create permission
			\item \texttt{SELECT} only
			\item \texttt{MODIFY} only
			\item \texttt{EXECUTE} only
		\end{enumerate}
	\end{frame}
	\begin{frame}{Answer 745}
		\textbf{Correct: A} --- registration needs catalog/schema use and create.\\
		\vspace{0.5em}
		\textbf{Wrong:}
		\begin{itemize}
			\item B, C, D are insufficient alone.
		\end{itemize}
	\end{frame}
	
	% ---------------------------------------------------------------
	\begin{frame}{Question 746}
		Which privilege is needed to serve a model from Unity Catalog?
		\begin{enumerate}[label=\Alph*.]
			\item \texttt{EXECUTE}
			\item \texttt{SELECT}
			\item \texttt{MODIFY}
			\item \texttt{READ}
		\end{enumerate}
	\end{frame}
	\begin{frame}{Answer 746}
		\textbf{Correct: A} --- \texttt{EXECUTE} allows model use.\\
		\vspace{0.5em}
		\textbf{Wrong:}
		\begin{itemize}
			\item B, C, D are table privileges.
		\end{itemize}
	\end{frame}
	
	% ---------------------------------------------------------------
	\begin{frame}{Question 747}
		What is a model version's URI in Unity Catalog?
		\begin{enumerate}[label=\Alph*.]
			\item \texttt{models:/catalog.schema.model/version}
			\item \texttt{runs:/version/model}
			\item \texttt{models:/catalog/version}
			\item \texttt{models://version}
		\end{enumerate}
	\end{frame}
	\begin{frame}{Answer 747}
		\textbf{Correct: A} --- UC model URIs use \texttt{models:/<name>/<version>}.\\
		\vspace{0.5em}
		\textbf{Wrong:}
		\begin{itemize}
			\item B is a run URI.
			\item C, D are malformed.
		\end{itemize}
	\end{frame}
	
	% ---------------------------------------------------------------
	\begin{frame}{Question 748}
		What does \texttt{mlflow.pyfunc.load\_model()} accept?
		\begin{enumerate}[label=\Alph*.]
			\item A model URI such as a registry URI or run URI
			\item A DataFrame
			\item A cluster
			\item A feature table
		\end{enumerate}
	\end{frame}
	\begin{frame}{Answer 748}
		\textbf{Correct: A} --- it loads by model URI.\\
		\vspace{0.5em}
		\textbf{Wrong:}
		\begin{itemize}
			\item B, C, D are not model URIs.
		\end{itemize}
	\end{frame}
	
	% ---------------------------------------------------------------
	\begin{frame}{Question 749}
		What does \texttt{mlflow.pyfunc.predict()} accept after loading?
		\begin{enumerate}[label=\Alph*.]
			\item A pandas or Spark DataFrame of inputs
			\item A model
			\item A feature table
			\item A cluster
		\end{enumerate}
	\end{frame}
	\begin{frame}{Answer 749}
		\textbf{Correct: A} --- \texttt{predict()} accepts input DataFrames.\\
		\vspace{0.5em}
		\textbf{Wrong:}
		\begin{itemize}
			\item B, C, D are not inputs.
		\end{itemize}
	\end{frame}
	
	% ---------------------------------------------------------------
	\begin{frame}{Question 750}
		What is an MLflow model URI \texttt{runs:/...} used for?
		\begin{enumerate}[label=\Alph*.]
			\item Loading a model artifact from a run
			\item Loading a feature table
			\item Loading a cluster
			\item Loading a notebook
		\end{enumerate}
	\end{frame}
	\begin{frame}{Answer 750}
		\textbf{Correct: A} --- run URIs load artifacts from a run.\\
		\vspace{0.5em}
		\textbf{Wrong:}
		\begin{itemize}
			\item B, C, D are not loaded by run URIs.
		\end{itemize}
	\end{frame}
	
	% ---------------------------------------------------------------
	\begin{frame}{Question 751}
		What is an MLflow run's source name?
		\begin{enumerate}[label=\Alph*.]
			\item The file or notebook that started the run
			\item The model
			\item The feature table
			\item The cluster
		\end{enumerate}
	\end{frame}
	\begin{frame}{Answer 751}
		\textbf{Correct: A} --- source name identifies the run origin.\\
		\vspace{0.5em}
		\textbf{Wrong:}
		\begin{itemize}
			\item B, C, D are unrelated.
		\end{itemize}
	\end{frame}
	
	% ---------------------------------------------------------------
	\begin{frame}{Question 752}
		What is an MLflow run's user?
		\begin{enumerate}[label=\Alph*.]
			\item The user who started the run
			\item The model
			\item The feature table
			\item The cluster
		\end{enumerate}
	\end{frame}
	\begin{frame}{Answer 752}
		\textbf{Correct: A} --- the run user is the initiator.\\
		\vspace{0.5em}
		\textbf{Wrong:}
		\begin{itemize}
			\item B, C, D are unrelated.
		\end{itemize}
	\end{frame}
	
	% ---------------------------------------------------------------
	\begin{frame}{Question 753}
		What is an MLflow run's start time?
		\begin{enumerate}[label=\Alph*.]
			\item When the run began
			\item When the run ended
			\item When the model was served
			\item When the model was registered
		\end{enumerate}
	\end{frame}
	\begin{frame}{Answer 753}
		\textbf{Correct: A} --- start time is when the run began.\\
		\vspace{0.5em}
		\textbf{Wrong:}
		\begin{itemize}
			\item B, C, D are other times.
		\end{itemize}
	\end{frame}
	
	% ---------------------------------------------------------------
	\begin{frame}{Question 754}
		What is an MLflow run's end time?
		\begin{enumerate}[label=\Alph*.]
			\item When the run ended
			\item When the run began
			\item When the model was served
			\item When the model was registered
		\end{enumerate}
	\end{frame}
	\begin{frame}{Answer 754}
		\textbf{Correct: A} --- end time is when the run ended.\\
		\vspace{0.5em}
		\textbf{Wrong:}
		\begin{itemize}
			\item B, C, D are other times.
		\end{itemize}
	\end{frame}
	
	% ---------------------------------------------------------------
	\begin{frame}{Question 755}
		What is an MLflow experiment's artifact location?
		\begin{enumerate}[label=\Alph*.]
			\item Where artifacts for runs in the experiment are stored
			\item Where models are registered
			\item Where clusters are configured
			\item Where notebooks are stored
		\end{enumerate}
	\end{frame}
	\begin{frame}{Answer 755}
		\textbf{Correct: A} --- artifact locations define storage for runs.\\
		\vspace{0.5em}
		\textbf{Wrong:}
		\begin{itemize}
			\item B, C, D are unrelated.
		\end{itemize}
	\end{frame}
	
	% ---------------------------------------------------------------
	\begin{frame}{Question 756}
		What is an MLflow experiment's lifecycle stage?
		\begin{enumerate}[label=\Alph*.]
			\item Whether the experiment is active or deleted
			\item The model's stage
			\item The cluster's mode
			\item The feature table's mode
		\end{enumerate}
	\end{frame}
	\begin{frame}{Answer 756}
		\textbf{Correct: A} --- lifecycle stages track active/deleted experiments.\\
		\vspace{0.5em}
		\textbf{Wrong:}
		\begin{itemize}
			\item B, C, D are unrelated.
		\end{itemize}
	\end{frame}
	
	% ---------------------------------------------------------------
	\begin{frame}{Question 757}
		What is a nested run in MLflow?
		\begin{enumerate}[label=\Alph*.]
			\item A run started inside another run
			\item A model
			\item A feature table
			\item A cluster
		\end{enumerate}
	\end{frame}
	\begin{frame}{Answer 757}
		\textbf{Correct: A} --- nested runs belong to a parent run.\\
		\vspace{0.5em}
		\textbf{Wrong:}
		\begin{itemize}
			\item B, C, D are unrelated.
		\end{itemize}
	\end{frame}
	
	% ---------------------------------------------------------------
	\begin{frame}{Question 758}
		Which flag starts a nested run?
		\begin{enumerate}[label=\Alph*.]
			\item \texttt{nested=True}
			\item \texttt{child=True}
			\item \texttt{sub=True}
			\item \texttt{nest=True}
		\end{enumerate}
	\end{frame}
	\begin{frame}{Answer 758}
		\textbf{Correct: A} --- \texttt{nested=True} starts a nested run.\\
		\vspace{0.5em}
		\textbf{Wrong:}
		\begin{itemize}
			\item B, C, D are not the API flag.
		\end{itemize}
	\end{frame}
	
	% ---------------------------------------------------------------
	\begin{frame}{Question 759}
		What is an MLflow parent run?
		\begin{enumerate}[label=\Alph*.]
			\item The outer run containing nested runs
			\item A model
			\item A feature table
			\item A cluster
		\end{enumerate}
	\end{frame}
	\begin{frame}{Answer 759}
		\textbf{Correct: A} --- the parent run contains nested runs.\\
		\vspace{0.5em}
		\textbf{Wrong:}
		\begin{itemize}
			\item B, C, D are unrelated.
		\end{itemize}
	\end{frame}
	
	% ---------------------------------------------------------------
	\begin{frame}{Question 760}
		What does \texttt{mlflow.get\_parent\_run()} return?
		\begin{enumerate}[label=\Alph*.]
			\item The parent run of the active nested run
			\item A model
			\item A feature table
			\item A cluster
		\end{enumerate}
	\end{frame}
	\begin{frame}{Answer 760}
		\textbf{Correct: A} --- it returns the parent run.\\
		\vspace{0.5em}
		\textbf{Wrong:}
		\begin{itemize}
			\item B, C, D are unrelated.
		\end{itemize}
	\end{frame}
	
	% ---------------------------------------------------------------
	\begin{frame}{Question 761}
		Which MLflow API lists nested runs?
		\begin{enumerate}[label=\Alph*.]
			\item \texttt{mlflow.search\_runs()} with a parent filter
			\item \texttt{mlflow.list\_children()}
			\item \texttt{mlflow.children()}
			\item \texttt{mlflow.nested\_runs()}
		\end{enumerate}
	\end{frame}
	\begin{frame}{Answer 761}
		\textbf{Correct: A} --- search with a parent run filter.\\
		\vspace{0.5em}
		\textbf{Wrong:}
		\begin{itemize}
			\item B, C, D are not the API.
		\end{itemize}
	\end{frame}
	
	% ---------------------------------------------------------------
	\begin{frame}{Question 762}
		What is a Feature Store training set's source?
		\begin{enumerate}[label=\Alph*.]
			\item A label DataFrame joined with feature lookups
			\item A model
			\item A cluster
			\item A notebook
		\end{enumerate}
	\end{frame}
	\begin{frame}{Answer 762}
		\textbf{Correct: A} --- training sets combine labels and features.\\
		\vspace{0.5em}
		\textbf{Wrong:}
		\begin{itemize}
			\item B, C, D are unrelated.
		\end{itemize}
	\end{frame}
	
	% ---------------------------------------------------------------
	\begin{frame}{Question 763}
		What does a training set's label column represent?
		\begin{enumerate}[label=\Alph*.]
			\item The target value for each row
			\item The feature table
			\item The model
			\item The cluster
		\end{enumerate}
	\end{frame}
	\begin{frame}{Answer 763}
		\textbf{Correct: A} --- label columns hold targets.\\
		\vspace{0.5em}
		\textbf{Wrong:}
		\begin{itemize}
			\item B, C, D are unrelated.
		\end{itemize}
	\end{frame}
	
	% ---------------------------------------------------------------
	\begin{frame}{Question 764}
		What does \texttt{exclude\_columns} do in a FeatureLookup?
		\begin{enumerate}[label=\Alph*.]
			\item Excludes specific columns from the lookup
			\item Trains a model
			\item Serves a model
			\item Registers a model
		\end{enumerate}
	\end{frame}
	\begin{frame}{Answer 764}
		\textbf{Correct: A} --- \texttt{exclude\_columns} removes columns.\\
		\vspace{0.5em}
		\textbf{Wrong:}
		\begin{itemize}
			\item B, C, D are unrelated.
		\end{itemize}
	\end{frame}
	
	% ---------------------------------------------------------------
	\begin{frame}{Question 765}
		What does a FeatureLookup's \texttt{lookup\_key} specify?
		\begin{enumerate}[label=\Alph*.]
			\item The join key used to match features to the label DataFrame
			\item The model
			\item The cluster
			\item The metric
		\end{enumerate}
	\end{frame}
	\begin{frame}{Answer 765}
		\textbf{Correct: A} --- \texttt{lookup\_key} is the join key.\\
		\vspace{0.5em}
		\textbf{Wrong:}
		\begin{itemize}
			\item B, C, D are unrelated.
		\end{itemize}
	\end{frame}
	
	% ---------------------------------------------------------------
	\begin{frame}{Question 766}
		What does a FeatureLookup's \texttt{feature\_names} specify?
		\begin{enumerate}[label=\Alph*.]
			\item The specific feature columns to look up
			\item The model
			\item The cluster
			\item The metric
		\end{enumerate}
	\end{frame}
	\begin{frame}{Answer 766}
		\textbf{Correct: A} --- \texttt{feature\_names} selects columns.\\
		\vspace{0.5em}
		\textbf{Wrong:}
		\begin{itemize}
			\item B, C, D are unrelated.
		\end{itemize}
	\end{frame}
	
	% ---------------------------------------------------------------
	\begin{frame}{Question 767}
		What is a model's input example used for in MLflow?
		\begin{enumerate}[label=\Alph*.]
			\item Documenting expected input format for serving
			\item Training a model
			\item Registering a model
			\item Deleting a model
		\end{enumerate}
	\end{frame}
	\begin{frame}{Answer 767}
		\textbf{Correct: A} --- input examples document expected inputs.\\
		\vspace{0.5em}
		\textbf{Wrong:}
		\begin{itemize}
			\item B, C, D are unrelated.
		\end{itemize}
	\end{frame}
	
	% ---------------------------------------------------------------
	\begin{frame}{Question 768}
		What does \texttt{mlflow.models.infer\_signature()} need to work?
		\begin{enumerate}[label=\Alph*.]
			\item Sample input data and a model
			\item A cluster
			\item A feature table
			\item A notebook
		\end{enumerate}
	\end{frame}
	\begin{frame}{Answer 768}
		\textbf{Correct: A} --- it needs sample input and model.\\
		\vspace{0.5em}
		\textbf{Wrong:}
		\begin{itemize}
			\item B, C, D are unrelated.
		\end{itemize}
	\end{frame}
	
	% ---------------------------------------------------------------
	\begin{frame}{Question 769}
		What is an MLflow model signature's input schema?
		\begin{enumerate}[label=\Alph*.]
			\item Column names and types expected at inference
			\item The model weights
			\item The training data
			\item The cluster
		\end{enumerate}
	\end{frame}
	\begin{frame}{Answer 769}
		\textbf{Correct: A} --- input schema defines expected columns/types.\\
		\vspace{0.5em}
		\textbf{Wrong:}
		\begin{itemize}
			\item B, C, D are unrelated.
		\end{itemize}
	\end{frame}
	
	% ---------------------------------------------------------------
	\begin{frame}{Question 770}
		What is an MLflow model signature's output schema?
		\begin{enumerate}[label=\Alph*.]
			\item Column names and types of predictions
			\item The model weights
			\item The training data
			\item The cluster
		\end{enumerate}
	\end{frame}
	\begin{frame}{Answer 770}
		\textbf{Correct: A} --- output schema defines prediction columns/types.\\
		\vspace{0.5em}
		\textbf{Wrong:}
		\begin{itemize}
			\item B, C, D are unrelated.
		\end{itemize}
	\end{frame}
	
	% ---------------------------------------------------------------
	\begin{frame}{Question 771}
		What does a model signature's \texttt{params} field describe?
		\begin{enumerate}[label=\Alph*.]
			\item Inference-time parameters
			\item Training data
			\item Model weights
			\item Cluster settings
		\end{enumerate}
	\end{frame}
	\begin{frame}{Answer 771}
		\textbf{Correct: A} --- \texttt{params} describes inference parameters.\\
		\vspace{0.5em}
		\textbf{Wrong:}
		\begin{itemize}
			\item B, C, D are unrelated.
		\end{itemize}
	\end{frame}
	
	% ---------------------------------------------------------------
	\begin{frame}{Question 772}
		Why is a model signature important for pipelines?
		\begin{enumerate}[label=\Alph*.]
			\item It catches input/output mismatches early
			\item It trains faster
			\item It serves faster
			\item It registers faster
		\end{enumerate}
	\end{frame}
	\begin{frame}{Answer 772}
		\textbf{Correct: A} --- signatures catch mismatches early.\\
		\vspace{0.5em}
		\textbf{Wrong:}
		\begin{itemize}
			\item B, C, D are false.
		\end{itemize}
	\end{frame}
	
	% ---------------------------------------------------------------
	\begin{frame}{Question 773}
		What is a Spark ML \texttt{OneHotEncoder} estimator output?
		\begin{enumerate}[label=\Alph*.]
			\item Sparse vectors
			\item Scalars
			\item Strings
			\item Buckets
		\end{enumerate}
	\end{frame}
	\begin{frame}{Answer 773}
		\textbf{Correct: A} --- one-hot outputs sparse vectors.\\
		\vspace{0.5em}
		\textbf{Wrong:}
		\begin{itemize}
			\item B, C, D are not outputs.
		\end{itemize}
	\end{frame}
	
	% ---------------------------------------------------------------
	\begin{frame}{Question 774}
		What is a Spark ML \texttt{StringIndexer} default behavior for unseen labels?
		\begin{enumerate}[label=\Alph*.]
			\item It can throw an error or assign a default index
			\item It ignores them
			\item It drops them
			\item It duplicates them
		\end{enumerate}
	\end{frame}
	\begin{frame}{Answer 774}
		\textbf{Correct: A} --- unseen labels can error or get defaults.\\
		\vspace{0.5em}
		\textbf{Wrong:}
		\begin{itemize}
			\item B, C, D are not its behavior.
		\end{itemize}
	\end{frame}
	
	% ---------------------------------------------------------------
	\begin{frame}{Question 775}
		What is a Spark ML \texttt{VectorIndexer} used for?
		\begin{enumerate}[label=\Alph*.]
			\item Indexing categorical columns inside a vector
			\item Scaling numeric features
			\item Imputing values
			\item Assembling vectors
		\end{enumerate}
	\end{frame}
	\begin{frame}{Answer 775}
		\textbf{Correct: A} --- VectorIndexer indexes categoricals in vectors.\\
		\vspace{0.5em}
		\textbf{Wrong:}
		\begin{itemize}
			\item B, C, D are other transformers.
		\end{itemize}
	\end{frame}
	
	% ---------------------------------------------------------------
	\begin{frame}{Question 776}
		What is a Spark ML \texttt{SQLTransformer} used for?
		\begin{enumerate}[label=\Alph*.]
			\item Applying a SQL query as a transformation
			\item Training a model
			\item Serving a model
			\item Registering a model
		\end{enumerate}
	\end{frame}
	\begin{frame}{Answer 776}
		\textbf{Correct: A} --- SQLTransformer applies SQL transformations.\\
		\vspace{0.5em}
		\textbf{Wrong:}
		\begin{itemize}
			\item B, C, D are unrelated.
		\end{itemize}
	\end{frame}
	
	% ---------------------------------------------------------------
	\begin{frame}{Question 777}
		What is a Spark ML \texttt{Imputer} default strategy?
		\begin{enumerate}[label=\Alph*.]
			\item Mean
			\item Median
			\item Mode
			\item Zero
		\end{enumerate}
	\end{frame}
	\begin{frame}{Answer 777}
		\textbf{Correct: A} --- the default strategy is mean.\\
		\vspace{0.5em}
		\textbf{Wrong:}
		\begin{itemize}
			\item B, C, D are other strategies.
		\end{itemize}
	\end{frame}
	
	% ---------------------------------------------------------------
	\begin{frame}{Question 778}
		What is a Spark ML \texttt{VectorAssembler} output?
		\begin{enumerate}[label=\Alph*.]
			\item A single feature vector column
			\item A scalar
			\item A string
			\item A bucket
		\end{enumerate}
	\end{frame}
	\begin{frame}{Answer 778}
		\textbf{Correct: A} --- it outputs a feature vector.\\
		\vspace{0.5em}
		\textbf{Wrong:}
		\begin{itemize}
			\item B, C, D are not outputs.
		\end{itemize}
	\end{frame}
	
	% ---------------------------------------------------------------
	\begin{frame}{Question 779}
		What is a Spark ML \texttt{Bucketizer} used for?
		\begin{enumerate}[label=\Alph*.]
			\item Binning continuous features
			\item Scaling features
			\item Encoding categories
			\item Imputing values
		\end{enumerate}
	\end{frame}
	\begin{frame}{Answer 779}
		\textbf{Correct: A} --- Bucketizer bins continuous features.\\
		\vspace{0.5em}
		\textbf{Wrong:}
		\begin{itemize}
			\item B, C, D are other transformers.
		\end{itemize}
	\end{frame}
	
	% ---------------------------------------------------------------
	\begin{frame}{Question 780}
		What is a Spark ML \texttt{QuantileDiscretizer} used for?
		\begin{enumerate}[label=\Alph*.]
			\item Binning continuous features by quantiles
			\item Scaling features
			\item Encoding categories
			\item Imputing values
		\end{enumerate}
	\end{frame}
	\begin{frame}{Answer 780}
		\textbf{Correct: A} --- quantile discretization uses quantiles.\\
		\vspace{0.5em}
		\textbf{Wrong:}
		\begin{itemize}
			\item B, C, D are other transformers.
		\end{itemize}
	\end{frame}
	
	% ---------------------------------------------------------------
	\begin{frame}{Question 781}
		What is a Spark ML \texttt{MinMaxScaler} output range?
		\begin{enumerate}[label=\Alph*.]
			\item $[0,1]$
			\item $[-1,1]$
			\item $[0,100]$
			\item Unbounded
		\end{enumerate}
	\end{frame}
	\begin{frame}{Answer 781}
		\textbf{Correct: A} --- MinMaxScaler outputs $[0,1]$ by default.\\
		\vspace{0.5em}
		\textbf{Wrong:}
		\begin{itemize}
			\item B, C, D are other ranges.
		\end{itemize}
	\end{frame}
	
	% ---------------------------------------------------------------
	\begin{frame}{Question 782}
		What is a Spark ML \texttt{StandardScaler} output?
		\begin{enumerate}[label=\Alph*.]
			\item Zero-mean, unit-variance features
			\item $[0,1]$ features
			\item $[-1,1]$ features
			\item Unbounded features
		\end{enumerate}
	\end{frame}
	\begin{frame}{Answer 782}
		\textbf{Correct: A} --- StandardScaler z-scores features.\\
		\vspace{0.5em}
		\textbf{Wrong:}
		\begin{itemize}
			\item B, C, D are other scales.
		\end{itemize}
	\end{frame}
	
	% ---------------------------------------------------------------
	\begin{frame}{Question 783}
		What is a Spark ML \texttt{Normalizer} used for?
		\begin{enumerate}[label=\Alph*.]
			\item Scaling each row to unit norm
			\item Scaling each column
			\item Encoding categories
			\item Imputing values
		\end{enumerate}
	\end{frame}
	\begin{frame}{Answer 783}
		\textbf{Correct: A} --- Normalizer scales rows to unit norm.\\
		\vspace{0.5em}
		\textbf{Wrong:}
		\begin{itemize}
			\item B, C, D are other transformers.
		\end{itemize}
	\end{frame}
	
	% ---------------------------------------------------------------
	\begin{frame}{Question 784}
		What is a Spark ML \texttt{MaxAbsScaler} used for?
		\begin{enumerate}[label=\Alph*.]
			\item Scaling by the maximum absolute value
			\item Scaling to $[0,1]$
			\item Encoding categories
			\item Imputing values
		\end{enumerate}
	\end{frame}
	\begin{frame}{Answer 784}
		\textbf{Correct: A} --- MaxAbsScaler divides by max abs.\\
		\vspace{0.5em}
		\textbf{Wrong:}
		\begin{itemize}
			\item B, C, D are other methods.
		\end{itemize}
	\end{frame}
	
	% ---------------------------------------------------------------
	\begin{frame}{Question 785}
		What is a Spark ML \texttt{PolynomialExpansion} output?
		\begin{enumerate}[label=\Alph*.]
			\item Polynomial features up to a degree
			\item A single feature
			\item A category
			\item A bucket
		\end{enumerate}
	\end{frame}
	\begin{frame}{Answer 785}
		\textbf{Correct: A} --- it outputs polynomial features.\\
		\vspace{0.5em}
		\textbf{Wrong:}
		\begin{itemize}
			\item B, C, D are not outputs.
		\end{itemize}
	\end{frame}
	
	% ---------------------------------------------------------------
	\begin{frame}{Question 786}
		What is a Spark ML \texttt{DCT} used for?
		\begin{enumerate}[label=\Alph*.]
			\item Discrete cosine transform on a vector
			\item Scaling
			\item Encoding
			\item Imputing
		\end{enumerate}
	\end{frame}
	\begin{frame}{Answer 786}
		\textbf{Correct: A} --- DCT applies a cosine transform.\\
		\vspace{0.5em}
		\textbf{Wrong:}
		\begin{itemize}
			\item B, C, D are other transformations.
		\end{itemize}
	\end{frame}
	
	% ---------------------------------------------------------------
	\begin{frame}{Question 787}
		What is a Spark ML \texttt{ElementwiseProduct} used for?
		\begin{enumerate}[label=\Alph*.]
			\item Multiplying each feature by a scalar
			\item Scaling to $[0,1]$
			\item Encoding categories
			\item Imputing values
		\end{enumerate}
	\end{frame}
	\begin{frame}{Answer 787}
		\textbf{Correct: A} --- ElementwiseProduct scales features.\\
		\vspace{0.5em}
		\textbf{Wrong:}
		\begin{itemize}
			\item B, C, D are other methods.
		\end{itemize}
	\end{frame}
	
	% ---------------------------------------------------------------
	\begin{frame}{Question 788}
		What is a Spark ML \texttt{Interaction} transformer used for?
		\begin{enumerate}[label=\Alph*.]
			\item Creating interaction features between vectors
			\item Scaling features
			\item Encoding categories
			\item Imputing values
		\end{enumerate}
	\end{frame}
	\begin{frame}{Answer 788}
		\textbf{Correct: A} --- Interaction creates interaction features.\\
		\vspace{0.5em}
		\textbf{Wrong:}
		\begin{itemize}
			\item B, C, D are other transformers.
		\end{itemize}
	\end{frame}
	
	% ---------------------------------------------------------------
	\begin{frame}{Question 789}
		What is a Spark ML \texttt{VectorSlicer} used for?
		\begin{enumerate}[label=\Alph*.]
			\item Selecting a subset of features from a vector
			\item Scaling features
			\item Encoding categories
			\item Imputing values
		\end{enumerate}
	\end{frame}
	\begin{frame}{Answer 789}
		\textbf{Correct: A} --- VectorSlicer selects features.\\
		\vspace{0.5em}
		\textbf{Wrong:}
		\begin{itemize}
			\item B, C, D are other transformers.
		\end{itemize}
	\end{frame}
	
	% ---------------------------------------------------------------
	\begin{frame}{Question 790}
		What is a Spark ML \texttt{ChiSqSelector} used for?
		\begin{enumerate}[label=\Alph*.]
			\item Selecting features using a chi-squared test
			\item Scaling features
			\item Encoding categories
			\item Imputing values
		\end{enumerate}
	\end{frame}
	\begin{frame}{Answer 790}
		\textbf{Correct: A} --- ChiSqSelector is a feature selector.\\
		\vspace{0.5em}
		\textbf{Wrong:}
		\begin{itemize}
			\item B, C, D are other transformers.
		\end{itemize}
	\end{frame}
	
	% ---------------------------------------------------------------
	\begin{frame}{Question 791}
		What is a Spark ML \texttt{VarianceThresholdSelector} used for?
		\begin{enumerate}[label=\Alph*.]
			\item Selecting features with variance above a threshold
			\item Scaling features
			\item Encoding categories
			\item Imputing values
		\end{enumerate}
	\end{frame}
	\begin{frame}{Answer 791}
		\textbf{Correct: A} --- it selects features by variance.\\
		\vspace{0.5em}
		\textbf{Wrong:}
		\begin{itemize}
			\item B, C, D are other transformers.
		\end{itemize}
	\end{frame}
	
	% ---------------------------------------------------------------
	\begin{frame}{Question 792}
		What is a Spark ML \texttt{UnivariateFeatureSelector} used for?
		\begin{enumerate}[label=\Alph*.]
			\item Selecting features via univariate statistical tests
			\item Scaling features
			\item Encoding categories
			\item Imputing values
		\end{enumerate}
	\end{frame}
	\begin{frame}{Answer 792}
		\textbf{Correct: A} --- it selects features statistically.\\
		\vspace{0.5em}
		\textbf{Wrong:}
		\begin{itemize}
			\item B, C, D are other transformers.
		\end{itemize}
	\end{frame}
	
	% ---------------------------------------------------------------
	\begin{frame}{Question 793}
		What is a Spark ML \texttt{RFormula} used for?
		\begin{enumerate}[label=\Alph*.]
			\item Specifying a model formula in R syntax
			\item Scaling features
			\item Encoding categories
			\item Imputing values
		\end{enumerate}
	\end{frame}
	\begin{frame}{Answer 793}
		\textbf{Correct: A} --- RFormula uses R-style formulas.\\
		\vspace{0.5em}
		\textbf{Wrong:}
		\begin{itemize}
			\item B, C, D are other transformers.
		\end{itemize}
	\end{frame}
	
	% ---------------------------------------------------------------
	\begin{frame}{Question 794}
		What is a Spark ML \texttt{DecisionTreeClassifier} used for?
		\begin{enumerate}[label=\Alph*.]
			\item Classification using decision trees
			\item Clustering
			\item Regression
			\item Dimensionality reduction
		\end{enumerate}
	\end{frame}
	\begin{frame}{Answer 794}
		\textbf{Correct: A} --- DecisionTreeClassifier is a classifier.\\
		\vspace{0.5em}
		\textbf{Wrong:}
		\begin{itemize}
			\item B, C, D are other tasks.
		\end{itemize}
	\end{frame}
	
	% ---------------------------------------------------------------
	\begin{frame}{Question 795}
		What is a Spark ML \texttt{DecisionTreeRegressor} used for?
		\begin{enumerate}[label=\Alph*.]
			\item Regression using decision trees
			\item Classification
			\item Clustering
			\item Dimensionality reduction
		\end{enumerate}
	\end{frame}
	\begin{frame}{Answer 795}
		\textbf{Correct: A} --- DecisionTreeRegressor is a regressor.\\
		\vspace{0.5em}
		\textbf{Wrong:}
		\begin{itemize}
			\item B, C, D are other tasks.
		\end{itemize}
	\end{frame}
	
	% ---------------------------------------------------------------
	\begin{frame}{Question 796}
		What is a Spark ML \texttt{GBTClassifier} used for?
		\begin{enumerate}[label=\Alph*.]
			\item Classification using gradient-boosted trees
			\item Clustering
			\item Regression
			\item Dimensionality reduction
		\end{enumerate}
	\end{frame}
	\begin{frame}{Answer 796}
		\textbf{Correct: A} --- GBTClassifier is a classifier.\\
		\vspace{0.5em}
		\textbf{Wrong:}
		\begin{itemize}
			\item B, C, D are other tasks.
		\end{itemize}
	\end{frame}
	
	% ---------------------------------------------------------------
	\begin{frame}{Question 797}
		What is a Spark ML \texttt{GBTRegressor} used for?
		\begin{enumerate}[label=\Alph*.]
			\item Regression using gradient-boosted trees
			\item Classification
			\item Clustering
			\item Dimensionality reduction
		\end{enumerate}
	\end{frame}
	\begin{frame}{Answer 797}
		\textbf{Correct: A} --- GBTRegressor is a regressor.\\
		\vspace{0.5em}
		\textbf{Wrong:}
		\begin{itemize}
			\item B, C, D are other tasks.
		\end{itemize}
	\end{frame}
	
	% ---------------------------------------------------------------
	\begin{frame}{Question 798}
		What is a Spark ML \texttt{RandomForestClassifier} used for?
		\begin{enumerate}[label=\Alph*.]
			\item Classification using random forests
			\item Clustering
			\item Regression
			\item Dimensionality reduction
		\end{enumerate}
	\end{frame}
	\begin{frame}{Answer 798}
		\textbf{Correct: A} --- RandomForestClassifier is a classifier.\\
		\vspace{0.5em}
		\textbf{Wrong:}
		\begin{itemize}
			\item B, C, D are other tasks.
		\end{itemize}
	\end{frame}
	
	% ---------------------------------------------------------------
	\begin{frame}{Question 799}
		What is a Spark ML \texttt{RandomForestRegressor} used for?
		\begin{enumerate}[label=\Alph*.]
			\item Regression using random forests
			\item Classification
			\item Clustering
			\item Dimensionality reduction
		\end{enumerate}
	\end{frame}
	\begin{frame}{Answer 799}
		\textbf{Correct: A} --- RandomForestRegressor is a regressor.\\
		\vspace{0.5em}
		\textbf{Wrong:}
		\begin{itemize}
			\item B, C, D are other tasks.
		\end{itemize}
	\end{frame}
	
	% ---------------------------------------------------------------
	\begin{frame}{Question 800}
		What is a Spark ML \texttt{LinearSVC} used for?
		\begin{enumerate}[label=\Alph*.]
			\item Linear support vector classification
			\item Clustering
			\item Regression
			\item Dimensionality reduction
		\end{enumerate}
	\end{frame}
	\begin{frame}{Answer 800}
		\textbf{Correct: A} --- LinearSVC is a linear classifier.\\
		\vspace{0.5em}
		\textbf{Wrong:}
		\begin{itemize}
			\item B, C, D are other tasks.
		\end{itemize}
	\end{frame}
	
\end{document}